\begin{tikzpicture}[x=0.92cm,y=0.78cm]
\path[use as bounding box] (-0.20,0.05) rectangle (17.05,8.05);

% Source quiver and its Type A_2 full subquiver.
\node[wqdiagramlabel] at (3.70,7.50) {$Q$};
% \node[wqdiagramlabel] at (1.70,7.20)
%   {$a_\tau^\vee=(a_1,a_2,a_3)$};
% \node[wqdiagramlabel] at (5.70,7.20)
%   {$a_{\tau'}=(b_1,b_2,b_3)$};
\node[wqweavefrozen] (source-tau) at (1.70,6.50) {};
\node[wqweavefrozen] (source-taup) at (5.70,6.50) {};
\node[wqvlabel,above left=0.8mm and 0.6mm of source-tau] {$\tau$};
\node[wqvlabel,above right=0.8mm and 0.6mm of source-taup] {$\tau'$};

\draw[wqqbox] (0.10,0.25) rectangle (7.30,3.55);
\node[wqqboxtitle] at (0.40,0.25) {$Q(iji)$};
\node[wqlevel,anchor=east] at (0.03,2.45) {$i$};
\node[wqlevel,anchor=east] at (0.03,0.95) {$j$};
\draw[wqguide] (0.17,2.45)--(7.17,2.45);
\draw[wqguide] (0.17,0.95)--(7.17,0.95);

\node[wqstringfrozen] (source-il) at (0.85,2.45)
  {$[0,1]$};
\node[wqpivot] (source-k) at (3.70,2.45)
  {$[1,3]$};
\node[wqstringfrozen] (source-ir) at (6.55,2.45)
  {$[3,\infty]$};
\node[wqstringfrozen] (source-jl) at (1.95,0.95)
  {$[0,2]$};
\node[wqstringfrozen] (source-jr) at (5.45,0.95)
  {$[2,\infty]$};

% All three edge families are superimposed.  Their planar crossings carry no
% additional data, so every arrow is drawn as one continuous path.
  % The three source triangles carrying a_tau^\vee.
  \wqextensiontriangle[][bend left=9]
    {source-il}{source-tau}{source-k}
  \wqextensiontriangle{source-jl}{source-tau}{source-jr}
  \wqextensiontriangle[bend left=9]
    {source-k}{source-tau}{source-ir}
  % The three source triangles carrying a_{tau'}.
  \wqextensiontriangle[][bend left=9]
    {source-il}{source-taup}{source-k}
  \wqextensiontriangle{source-jl}{source-taup}{source-jr}
  \wqextensiontriangle[bend left=9]
    {source-k}{source-taup}{source-ir}
  % The string quiver and the local weave-quiver edge.
  \draw[wqstringarrow] (source-k)--(source-il);
  \draw[wqstringarrow] (source-ir)--(source-k);
  \draw[wqstringarrow] (source-jr)--(source-jl);
  \draw[wqstringhalf] (source-il)--(source-jl);
  \draw[wqstringarrow] (source-jl)--(source-k);
  \draw[wqstringarrow] (source-k)--(source-jr);
  \draw[wqstringhalf] (source-jr)--(source-ir);
  \draw[wqlocalarrow] (source-tau)--(source-taup);

% Edge labels are a separate semantic layer.
\path (source-tau)--node[wqdiagramlabel,above=0.8mm] {$x$}
  (source-taup);
\path (source-il.north)--node[wqweightlabel,pos=0.28]
  {$a_1$}(source-tau);
\path (source-tau) to[bend left=9]
  node[wqweightlabel,pos=0.28]
  {$a_1$}(source-k.north);
\path (source-jl.north)--node[wqweightlabel,pos=0.28]
  {$a_2$}(source-tau);
\path (source-tau)--node[wqweightlabel,pos=0.28]
  {$a_2$}(source-jr.north);
\path (source-k.north) to[bend left=9]
  node[wqweightlabel,pos=0.28]
  {$a_3$}(source-tau);
\path (source-tau)--node[wqweightlabel,pos=0.28]
  {$a_3$}(source-ir.north);
\path (source-il.north)--node[wqweightlabel,pos=0.28]
  {$b_1$}(source-taup);
\path (source-taup) to[bend left=9]
  node[wqweightlabel,pos=0.28]
  {$b_1$}(source-k.north);
\path (source-jl.north)--node[wqweightlabel,pos=0.28]
  {$b_2$}(source-taup);
\path (source-taup)--node[wqweightlabel,pos=0.28]
  {$b_2$}(source-jr.north);
\path (source-k.north) to[bend left=9]
  node[wqweightlabel,pos=0.28]
  {$b_3$}(source-taup);
\path (source-taup)--node[wqweightlabel,pos=0.28]
  {$b_3$}(source-ir.north);
% \node[wqdiagramlabel,below=0.7mm of source-k] {$k$};

% The local cluster transformation and its relabeling.
\draw[wqtimearrow] (7.72,2.25)--(8.43,2.25);
\node[wqmovelabel,above=0.9mm] at (8.075,2.25)
  {$\mu_\sigma$};

% Target quiver and its relabeled Type A_2 full subquiver.
\node[wqdiagramlabel] at (12.90,7.50) {$Q'$};
% \node[wqdiagramlabel] at (10.80,7.20)
%   {$(a'_\tau)^\vee=\Phi_\sigma^\vee(a_\tau^\vee)$};
% \node[wqdiagramlabel] at (15.00,7.20)
%   {$a'_{\tau'}=\Phi_\sigma(a_{\tau'})$};
\node[wqweavefrozen] (target-tau) at (9.95,6.50) {};
\node[wqweavefrozen] (target-taup) at (15.85,6.50) {};
\node[wqvlabel,above left=0.8mm and 0.6mm of target-tau] {$\tau$};
\node[wqvlabel,above right=0.8mm and 0.6mm of target-taup] {$\tau'$};

\draw[wqqbox] (8.95,0.25) rectangle (16.85,3.55);
\node[wqqboxtitle] at (9.25,0.25) {$Q(jij)$};
\node[wqlevel,anchor=east] at (8.88,2.45) {$i$};
\node[wqlevel,anchor=east] at (8.88,0.95) {$j$};
\draw[wqguide] (9.02,2.45)--(16.72,2.45);
\draw[wqguide] (9.02,0.95)--(16.72,0.95);

\node[wqstringfrozen] (target-il) at (10.55,2.45)
  {$[0,2]$};
\node[wqstringfrozen] (target-ir) at (15.25,2.45)
  {$[2,\infty]$};
\node[wqstringfrozen] (target-jl) at (9.55,0.95)
  {$[0,1]$};
\node[wqpivot] (target-k) at (12.90,0.95)
  {$[1,3]$};
\node[wqstringfrozen] (target-jr) at (16.25,0.95)
  {$[3,\infty]$};

% The target uses the same continuous-edge convention as the source.
  % The three target triangles carrying (a'_tau)^\vee.
  \wqextensiontriangle[][bend left=9]
    {target-jl}{target-tau}{target-k}
  \wqextensiontriangle{target-il}{target-tau}{target-ir}
  \wqextensiontriangle[bend left=9]
    {target-k}{target-tau}{target-jr}
  % The three target triangles carrying a'_{tau'}.
  \wqextensiontriangle[][bend left=9]
    {target-jl}{target-taup}{target-k}
  \wqextensiontriangle{target-il}{target-taup}{target-ir}
  \wqextensiontriangle[bend left=9]
    {target-k}{target-taup}{target-jr}
  % The string quiver and the local weave-quiver edge.
  \draw[wqstringarrow] (target-ir)--(target-il);
  \draw[wqstringarrow] (target-k)--(target-jl);
  \draw[wqstringarrow] (target-jr)--(target-k);
  \draw[wqstringhalf] (target-jl)--(target-il);
  \draw[wqstringarrow] (target-il)--(target-k);
  \draw[wqstringarrow] (target-k)--(target-ir);
  \draw[wqstringhalf] (target-ir)--(target-jr);
  \draw[wqlocalarrow] (target-tau)--(target-taup);

% Edge labels are kept above the edge system.
\path (target-tau)--node[wqdiagramlabel,above=0.8mm]
  {$x+\Omega_{\mathbf b'}(a',b')-\Omega_{\mathbf b}(a,b)$}
  (target-taup);
\path (target-jl.north)--node[wqweightlabel,pos=0.28]
  {$a'_1$}(target-tau);
\path (target-tau) to[bend left=9]
  node[wqweightlabel,pos=0.28]
  {$a'_1$}(target-k.north);
\path (target-il.north)--node[wqweightlabel,pos=0.28]
  {$a'_2$}(target-tau);
\path (target-tau)--node[wqweightlabel,pos=0.28]
  {$a'_2$}(target-ir.north);
\path (target-k.north) to[bend left=9]
  node[wqweightlabel,pos=0.28]
  {$a'_3$}(target-tau);
\path (target-tau)--node[wqweightlabel,pos=0.28]
  {$a'_3$}(target-jr.north);
\path (target-jl.north)--node[wqweightlabel,pos=0.28]
  {$b'_1$}(target-taup);
\path (target-taup) to[bend left=9]
  node[wqweightlabel,pos=0.28]
  {$b'_1$}(target-k.north);
\path (target-il.north)--node[wqweightlabel,pos=0.28]
  {$b'_2$}(target-taup);
\path (target-taup)--node[wqweightlabel,pos=0.28]
  {$b'_2$}(target-ir.north);
\path (target-k.north) to[bend left=9]
  node[wqweightlabel,pos=0.28]
  {$b'_3$}(target-taup);
\path (target-taup)--node[wqweightlabel,pos=0.28]
  {$b'_3$}(target-jr.north);
% \node[wqdiagramlabel,below=0.7mm of target-k] {$k'$};
\end{tikzpicture}
